\section{Evaluation}
\label{sec:evaluation}

\subsection{Vermont: Zero-Variance Baseline}
\label{sec:eval:vt}

Vermont has one congressional district (the entire state). Every seed produces
the identical assignment --- the trivial partition assigning all 192 census
tracts to district 1. The edge-cut is zero, the proportionality gap is
$100(1.0 - 0.683) = +31.7$\,pp regardless of seed, and the assignment files
are byte-identical across all 20 tested seeds. This confirms that for
at-large states, the seed is \emph{mathematically irrelevant}: no bisection
is performed and no partitioning choice is made.

\subsection{Pennsylvania: The Dense Solution Space}
\label{sec:eval:pa}

Pennsylvania (17 districts, 50.6\% Dem statewide vote, 2020 presidential)
is the paper's central case. The achievable proportionality gaps are
$\{-15.3, -9.4, -3.5, +2.3\}$\,pp corresponding to 6D, 7D, 8D, and 9D
Republican/Democratic seat splits.

\paragraph{Solution-space density.}
Across 100 tested seeds, every seed produced a \emph{distinct} final
partition (100 unique SHA-256 fingerprints of the assignment files). The
solution space is dense: there is no small set of ``canonical'' boundaries
that seeds converge to. This is a direct consequence of the high dimensionality
of the partition space ($\sim$193 census tracts, $\binom{193}{k}$-scale
feasible-partition count).

\paragraph{Edge-cut distribution.}
The observed edge-cuts span 2{,}529--3{,}003\,km (18.7\% range). Edge-cut
distributions overlap substantially across seat-outcome classes:

\begin{table}[h]
\centering
\begin{tabular}{lrrrr}
\toprule
Outcome & $n$ (100 seeds) & EC min & EC mean & EC max \\
\midrule
6D/11R ($-15.3$\,pp) & 9  & 2{,}616 & 2{,}720 & 2{,}892 \\
7D/10R ($-9.4$\,pp)  & 54 & 2{,}541 & 2{,}718 & 3{,}003 \\
8D/9R  ($-3.5$\,pp)  & 35 & 2{,}529 & 2{,}696 & 2{,}894 \\
9D/8R  ($+2.3$\,pp)  & 2  & 2{,}536 & 2{,}679 & 2{,}822 \\
\bottomrule
\end{tabular}
\caption{Edge-cut (km) by seat outcome, Pennsylvania, 100 seeds.
  All four outcome classes contain plans with edge-cuts within 5\% of
  the global minimum, confirming that minimum-edge-cut selection does not
  deterministically favour any partisan outcome.}
\label{tab:pa-ec-by-outcome}
\end{table}

\paragraph{Convergence.}
The running minimum improved 7 times across seeds 1--100 (at seeds 1, 2, 3,
5, 12, 18, 34; Table~\ref{tab:pa-convergence}). The final 66 seeds produced
no improvement. The current best plan (seed 34, EC = 2{,}529\,km) produces
8D/9R ($-3.5$\,pp). % TODO: update with full 1000-seed results

\begin{table}[h]
\centering
\begin{tabular}{rrrl}
\toprule
Seed & New minimum (km) & Improvement & Outcome \\
\midrule
1   & 3{,}003 & --- (first) & 7D \\
2   & 2{,}822 & $-$181     & 9D \\
3   & 2{,}632 & $-$190     & 8D \\
5   & 2{,}606 & $-$26      & 8D \\
12  & 2{,}596 & $-$10      & 7D \\
18  & 2{,}578 & $-$18      & 8D \\
34  & 2{,}529 & $-$49      & 8D \\
% TODO: insert seeds 101-1000 improvements here
\bottomrule
\end{tabular}
\caption{Running minimum edge-cut improvements, Pennsylvania, seeds 1--100.
  % TODO: extend to 1000 seeds.}
\label{tab:pa-convergence}
\end{table}

\paragraph{Near-minimum partisan distribution.}
Of the 29 plans within 5\% of the current global minimum (EC $\leq$ 2{,}655\,km):
2 produce 6D, 15 produce 7D, 11 produce 8D, and 1 produces 9D. The near-minimum
region of the solution space spans all four outcome classes. This means:
\emph{the minimum edge-cut criterion does not deterministically select a
particular partisan outcome in Pennsylvania}. The outcome of the MEC plan
depends on which local optimum is the global minimum --- a question only
resolved by exhaustive (or near-exhaustive) seed sampling.

\subsection{Texas: Large-State Failures and Convergence Rate}
\label{sec:eval:tx}

Texas (38 districts) exhibited a 4\% seed-failure rate (seeds 1, 2, 6, 16
out of 100 produced population-balance violations exceeding the 0.5\%
tolerance). These failures are expected: with 38 districts and $\lceil
\log_2 38 \rceil = 6$ levels of recursion, small per-level imbalances
compound. The statute's tiered tolerance schedule (\S~104(b)(4)) addresses
this; our implementation applies a flat 0.5\% tolerance. % TODO: implement
% tiered tolerance and re-run.

Among 96 successful seeds, edge-cuts ranged from [TODO] to [TODO]\,km with
[TODO] distinct seat outcomes observed. Convergence is expected to be slower
than Pennsylvania due to the larger partition space.

\subsection{50-State Summary}
\label{sec:eval:national}

% TODO: populate from Phase 2 sweep (50 states × 1000 seeds)

Key hypotheses to evaluate:
\begin{itemize}
  \item \textbf{H1 (at-large states)}: States with $k=1$ (VT, AK, WY, DE, ND,
        SD, MT) have zero seed variance. \emph{Confirmed for Vermont.}
  \item \textbf{H2 (competitive states)}: States with 10--20 competitive
        districts show the largest seed variance in partisan outcomes.
        \emph{Pennsylvania preliminary data: 3-seat swing across 100 seeds.}
  \item \textbf{H3 (sorted states)}: Deep-red and deep-blue states show low
        seed variance because geography dominates. \emph{TBD.}
  \item \textbf{H4 (MEC stability)}: The minimum-edge-cut plan is a typical
        draw from the near-minimum distribution, not a partisan outlier.
        \emph{Pennsylvania: MEC plan (8D) sits in the interior of the
        distribution, not at the partisan extreme.}
\end{itemize}
